\documentclass{article}
\usepackage[letterpaper, margin=1in]{geometry}
\usepackage{amsmath,amsthm,amssymb,mathtools,bm}
\usepackage{tikz-cd}
\usepackage{graphicx}
\usepackage{xcolor}

\newtheoremstyle{spacious}% <name>
  {8pt}      % Space above
  {16pt}     % Space below  <-- increase this for more after-spacing
  {\normalfont} % Body font
  {}         % Indent amount
  {\bfseries}% Head font
  {.}        % Punctuation after head
  {0.5em}    % Space after head
  {}         % Head spec

\theoremstyle{spacious}
\newtheorem*{dfn}{Dfn}
\newtheorem*{prp}{Proposition}
\newtheorem*{thm}{Theorem}
\newtheorem*{rmk}{Remark}
\newtheorem*{lem}{Lemma}
\newtheorem*{cor}{Corollary}
\newtheorem*{stp}{Step Result}
\newtheorem*{prm}{Exercise}
\graphicspath{ {./images/} }

\newcommand{\vl}{\vartriangleleft}
\newcommand{\vr}{\vartriangleright}

\newenvironment{pf}{\begin{enumerate} \item[] \textit{\underline{Proof}} }{\end{enumerate}\qed}
\renewcommand{\thefootnote}{\fnsymbol{footnote}}
\setcounter{footnote}{1}  
\renewcommand{\:}{\vcentcolon=} % for pretty :=
\newcommand{\Q}{\mathbb{Q}}
\newcommand{\Z}{\mathbb{Z}}
\newcommand{\R}{\mathbb{R}}
\newcommand{\N}{\mathbb{N}}
\newcommand{\C}{\mathbb{C}}
\newcommand{\F}{\mathcal{F}}
\renewcommand{\k}{\Bbbk}
\newcommand{\st}{\quad \text{s.t} \quad}

\newcommand{\just}[2]{\underset{\rotatebox{90}{\tiny\text{#2 }}}{#1}}
\newcommand{\jeq}[1]{\underset{\rotatebox{90}{\tiny\text{ #1}}}{=}}

\DeclareMathOperator{\Hom}{Hom}

% \DeclarePairedDelimiter{\ceil}{\lceil}{\rceil}
% \DeclarePairedDelimiter{\floor}{\lfloor}{\rfloor}

% Category Specific
\newcommand{\epi}{\twoheadrightarrow}
\newcommand{\mono}{\rightarrowtail}
\newcommand{\op}{\mathrm{op}}

\renewcommand{\labelenumi}{(\textit{\roman{enumi}})}
\renewcommand{\theenumi}{\roman{enumi}} % what \ref prints

\usepackage{subfiles} % Best loaded last in the preamble

\begin{document}

% \section*{0 tikzcd Notes}

\section*{1.1 Abstract and concrete categories}

\begin{dfn} A \textbf{category} consists of 
    \begin{itemize}
        \item a collection of \textbf{objects} \(X,Y,Z,\ldots\)
        \item a collection of \textbf{morphisms} \(f,g,h,\ldots\)
    \end{itemize}
So that:
\begin{itemize}
    \item Each morphism has specified \textbf{domain} and \textbf{codomain} objects; the notation \(f:X \to Y\) signifies that \(f\) is a morphism with domain \(X\) and codomain \(Y\)
    \item Each object has a designated \textbf{identity morphism} \(1_X:X \to X\)
    \item For any pair of morphism \(f,g\)  with the codomain of f equal to the domain of \(g\), there exists a specified \textbf{composite morphism} \(g f\) whose domain is equal to the domain of f and whose codomain is equal to the codomain of \(g\), i.e.,:
    \[
        f: X\to Y,\qquad g:Y\to Z \quad \leadsto \quad gf: X \to Z
    \]
\end{itemize}
This data is subject to the following two axioms:
\begin{itemize}
    \item For any \(f:X \to Y\), the composites \(1_Y f\) and \(f 1_X\) are both equal to \(f\).
    \item For any composable triple of morphisms \(f,g,h,\) the composites \(h(gf)\) and \((hg)f\) are equal and henceforth denoted by \(hgf\)
    \[
        f:X\to Y, \quad g:Y\to Z, \quad h:Z\to W \quad \leadsto \quad hgf: X \to W
    \]
That is, the composition law is associative and unital with the identity morphisms serving as two-sided identites.
\end{itemize}
\end{dfn}

\begin{dfn} A category is \textbf{concrete} when its objects are actually described as sets and its morphisms are actually functions between them and we refer to such morphisms as maps.  Otherwise, the category is \textbf{abstract} and we refer to the morphisms as arrows
    
\end{dfn}

\begin{dfn} A category is \textbf{small} if it has only a set's worth of morphisms.
\end{dfn}

\begin{dfn} A category is \textbf{locally small} if between any pair of objects there is only a set's worth of morphisms.  Then we \textit{are justified} in writing:
\[
    \mathsf{C}(X,Y) \quad \text{or} \quad \mathrm{Hom}\left(X,Y\right)
\]
and we call it a \textbf{hom-set}, but we \textit{abusively} use the same notation for a \textbf{hom-class}.
\end{dfn}

\begin{dfn} An \textbf{isomorphism} in a category is a morphism \(f:X\to Y\) for which there exists a morphism \(g:Y\to X\) so that \(gf=1_X\) and $fg=1_Y$.  The objects \(X\) and \(Y\) are \textbf{isomorphic} whenever there exists an isomorphism between \(X\) and \(Y\), in which case one writes \(X \cong Y\).

An \textbf{endomorphism} is  a morphism with equal domain and codomain

An \textbf{automorphism} is an isomorphism that is an endomorphism

\end{dfn}

\begin{dfn} A \textbf{groupoid} is a category in which every morphism is an isomorphism.
\end{dfn}

\begin{dfn} A \textbf{subcategory} \(\mathsf{D}\) of a category \(\mathsf{C}\) is a restriction of a category to a subcollection of objects and a subcollection of morphisms such that 
    \begin{enumerate}
        \item The domains and codomains of morphisms in \(\mathsf{D}\) are contained in \(\mathsf{D}\)
        \item The identities of all objects in \(\mathsf{D}\) are contained in \(\mathsf{D}\)
        \item The compositions of every composable pair of morphisms in \(\mathsf{D}\) are contained in \(\mathsf{D}\)
    \end{enumerate}
\end{dfn}

\begin{lem} Any category $\mathsf{C}$ contains a \textbf{maximal groupoid}, the subcategory containing all of the objects and only those morphisms that are isomorphisms.    
\end{lem}

\begin{prm}[1.1.i]
\phantom{t}
    \begin{enumerate}
        \item consider a morphism \(f:x\to y\).  Show that if there exists a pair of morphisms \(g,h: y \rightrightarrows x\) so that \(gf=1_x\) and \(fh=1_y\), then \(g=h\) and \(f\) is an isomorphism.
        \item Show that a morphism can have at most one inverse isomorphism.
    \end{enumerate}
\end{prm}

\begin{pf}{(i)}
    \item[] Assume \(\exists g,h:y\rightrightarrows x\) so that \(gf=1_x\) and \(fh=1_y\).  Then
    \[
        g\jeq{morphism dfn.} g1_y\jeq{asumptions }g(fh)\jeq{associativity }(gf)h\jeq{assumptions}1_xh\jeq{identity}h
    \]
\end{pf}

\begin{pf}{(ii)}
    \item[] Let \(g,h:y\rightrightarrows x\) be inverse isomorphisms of \(f:x\to y\)
    \item[] But this implies \(gf=1_x\) and \(fh=1_y\)
    \item[] Hence by the previous arguement \(g=h\)
\end{pf}

\begin{prm}[1.1.ii] Let \(\mathsf{C}\) be a category.  Show that the collection of isomorphisms in \(\mathsf{C}\) defines a subcategory, the \textbf{maximal groupoid} inside \(\mathsf{C}\)
\end{prm}

\begin{pf}{}
    \item[] Let \(\mathsf{C}\) be a category
    \item[] Consider the following data \(\mathsf{D}\) (we haven't yet confirmed it is a category): 
    \begin{enumerate}
    \item[] \(\mathrm{Obj}\left(\mathsf{D}\right):=\mathrm{Obj}\left(\mathsf{C}\right)\)
    \item[] \(x,y\in\mathrm{Obj}\left(\mathsf{D}\right)\bm{\implies}\mathsf{D}\left(x,y\right):=\{f \in \mathsf{C}\left(x,y\right) \mid f \text{ is an isomorphism}\}\)
    \end{enumerate}
    
    \item[] It suffices to verify the $3$ subcategory conditions for \(\mathsf{D}\):
    \begin{enumerate}
    \item[(i)] All objects are still present so trivially the domains and codomains of all the \textit{iso}morphisms are still present
    \item[(ii)] Since identities are always automorphisms they remain present in \(\mathsf{D}\)
    \item[(iii)] The composite of compossible pairs of isomorphisms are themselves isomorphisms \((f^{-1}g^{-1})(gf)=1_X\), hence are contained in \(\mathsf{D}\)
    \end{enumerate}
\end{pf}

\begin{prm}[1.1.iii.(i)]
    For any category \(\mathsf{C}\) and any object \(c\in \mathsf{C}\) there is a category \(c/ \mathsf{C}\) whose objects are morphisms \(f: c\to x\) with domain \(c\) and in which a morphism from \(f:c \to x\) to \(g:c \to y\) is a map \(h:x\to y\) between the codomains so that the triangle
        \[
        \begin{tikzcd}[column sep=small]
                        & c \ar[dl,"f"'] \ar[dr,"g"]     & \\
            x \ar[rr,"h"']   &                       & y
        \end{tikzcd}
        \]
        \textbf{commutes} (i.e. so that \(g=hf\))
\end{prm}

\begin{pf}{(i)}
    \item[] Let \(\mathsf{C}\) be a category with \(c\in\mathrm{Obj}\left(\mathsf{C}\right)\)
    \item[] Consider the following data \(c/\mathsf{C}\) (we haven't yet confirmed it is a category):
    \begin{enumerate}
        \item[] \(\mathrm{Obj}\left(c/\mathsf{C}\right):= \left\{f : c \to x \mid x \in \mathrm{Obj}(\mathsf{C})\right\}\) 
        \item[] \(f:c\to x,g:c\to y\in\mathrm{Obj}\left(c/\mathsf{C}\right)\bm{\implies}(\Hom_{c/\mathsf{C}})\left(f,g\right):=\{h : x \to y \in \mathsf{C}\mid g = h\circ f\}\)
        % \item[] \(\mathrm{Mor}\left(\mathsf{D}\right):= \{(f)\in \mathrm{Mor}\left(\mathsf{C}\right)\mid \mathrm{dom}\left(f\right)\in \mathrm{Obj}\left(\mathsf{D}\right) \land \mathrm{cod}\left(f\right)\in \mathrm{Obj}\left(\mathsf{D}\right) \}\)
    
    \end{enumerate}
            \item[] left composition with \(1_f\) gives the identity of \(f\in \mathrm{Mor}\left(\mathsf{D}\right)\)
        \item[] for all \(f,g\in\mathrm{Mor}\left(D\right)\) we can define \(fg\) as \(f\circ g \circ \)
        \item[] transitivity must hold since it holds for \(\mathsf{C}\)
\end{pf}


\begin{prm}[1.1.iii.(ii)]
    For any category \(\mathsf{C}\) and any object \(c\in \mathsf{C}\) there is a category \(\mathsf{C}/c\) whose objects are morphisms \(f: x\to c\) with domain \(c\) and in which a morphism from \(f:x \to c\) to \(g:y \to c\) is a map \(h:x\to y\) between the domains so that \(f=gh\).
\end{prm}



\begin{pf}{(ii)}
    \item[] Essentially the same arguement but using right composition
\end{pf}

\newpage

\section*{1.2. Duality}

\begin{dfn} Let \(\mathsf{C}\) be a category.  The \textbf{opposite category} \(\mathsf{C^{op}}\) consists of 
    \begin{itemize}
        \item The objects of \(\mathsf{C}\)
        \item For each morphisms $f:X \to Y$ in \(\mathsf{C}\) a morphism \(f^{\mathsf{op}}:Y \to X\) with domain and codomain swapped.
    \end{itemize}
\end{dfn}

\begin{dfn}[Post/Pre-composition] Fix \(\mathsf{C}\) with \(f:x \to y\) a morphism in \(\mathsf{C}\).
\begin{enumerate}
    \item define the post-composition of $f$:
    \[
        \forall c \in \mathsf{C} \quad f_*: \mathsf{C}(c,x) \to \mathsf{C}(c,y) \text{ by the rule } f_*(u):=u\mapsto f\circ u
    \]
    \item and define the pre-composition of $f$:
    \[
        \forall c \in \mathsf{C}\quad f^*: \mathsf{C}(y,c) \to \mathsf{C}(x,c) \text{ by the rule } f_*(v):=v\mapsto v\circ f
    \]
\end{enumerate}
\end{dfn}

\begin{lem}[Lemma 1.2.3.i] \(\forall c \in \mathsf{C}\)
    \[
        f:x\to y \text{ is an isomorphism in }\mathsf{C} \iff f_* \text{ is an isomorphism in \(\mathsf{Set}\)}
    \]
\end{lem}

\begin{rmk}
    \textit{put plainly we are saying that if there is an isomorphism between two objects in a category then if we consider the \underline{pair of sets} of morphisms \textbf{to} (and \textbf{from}) the two objects to any particular fixed object then those morphisms can be put into a bijection.}
\end{rmk}

\begin{pf}\(\left(\Rightarrow\right)\)
    \item[] Assume \(f:x\to y \text{ is an isomorphism in }\mathsf{C} \).
    \item[] Then, it has inverse \(g:y \to x\).
    \item[] Consider the post-composition of \(u\in\mathsf{C}\left(c,y\right)\) with $g$:
    \[
        g_*:\mathsf{C} \left(c,y\right) \to \mathsf{C}\left(c,x\right) \qquad g_*:=u\mapsto g\circ u
    \]
\item[] Take \(h \in \mathsf{C}(c,x)\)
\[\begin{aligned}
g_*f_*(h)&=g_*(f(h)) \quad \forall f \in \mathsf{C}\left(x,y\right)\quad &\text{dfn. post-comp.}\\
&=g(f(h)) \quad \phantom{{}_*}\forall f \in \mathsf{C}\left(x,y\right)\quad\forall g \in \mathsf{C}\left(y,x\right)&\text{dfn. post-comp.}\\
&=(gf)h \quad \phantom{{}_*}\phantom{()}\forall f \in \mathsf{C}\left(x,y\right)\quad\forall g \in \mathsf{C}\left(y,x\right)&\text{assoc}\\
&=1_x h \quad \phantom{{}_*}&\text{dfn. of inv. func.}\\
&=h \quad \phantom{{}_*}&\text{identity}\\
\end{aligned}\]
\item[] since this is true for arbitrary \(h\in \mathsf{C}(c,x)\) we see \(g_*f_*=1_{\mathsf{C}\left(c,x\right)}\)
\item[] Take \(k \in \mathsf{C}(c,y)\)
\[\begin{aligned}
f_*g_*(k)&=f_*g(k) \phantom{()}\quad \forall g \in \mathsf{C}\left(y,x\right)\quad &\text{dfn. post-comp.}\\
&=f(g(k)) \quad \phantom{{}_*}\forall g \in \mathsf{C}\left(y,x\right)\quad \forall f \in \mathsf{C}\left(x,y\right)&\text{dfn. post-comp.}\\
&=(fg)k \quad \phantom{{}_*}\phantom{()}\forall g \in \mathsf{C}\left(y,x\right)\quad \forall f \in \mathsf{C}\left(x,y\right)&\text{assoc}\\
&=1_y k \quad \phantom{{}_*}&\text{dfn. of inv. func.}\\
&=k \quad \phantom{{}_*}&\text{identity}\\
\end{aligned}\]
\item[] since this is true for arbitrary \(k\in \mathsf{C}(c,y)\) we see \(f_*g_*=1_{\mathsf{C}\left(c,y\right)}\)
\item[] Hence, since \(g_*f_*=1_{\mathsf{C}\left(c,x\right)}\) and \(f_*g_*=1_{\mathsf{C}\left(c,y\right)}\) we see that \(g_*\) is the inverse morphism of \(f_*\) in the category of \(\mathsf{Set}\)
\end{pf}

\begin{pf} \(\left(\Leftarrow\right)\)
    \item[] Assume  \(f_*:\mathsf{C}(c,x) \to \mathsf{C}(c,y)\) is an isomorphism in \(\mathsf{Set}\) and call its inverse morphism $\mathfrak{g}$:
    \[
    \mathfrak{g}f_*=1_{\mathsf{C}\left(c,x\right)} \text{ and } f_*\mathfrak{g}=1_{\mathsf{C}\left(c,y\right)}
    \]  
    {\textcolor{red}{\textbf{MUST COMEBACK AND FINISH}}}
\end{pf}

\begin{lem}[1.2.3.ii]
    \[
        f_* \text{ has an inverse} \leftrightsquigarrow f^* \text{ has an inverse}
    \]
\end{lem}

\begin{pf}{}
    \item[] {\textcolor{red}{\textbf{MUST COMEBACK AND FINISH}}}
\end{pf}

\begin{dfn} Let \(\mathsf{C}\) be a category with morphism \(f:x \to y\).  Then $f$ is...
    \begin{enumerate}
        \item a \textbf{monomorphism} (or just a \textbf{mono}) if for any parallel morphisms $h,k:w \rightrightarrows x$
        \[
            fh=fk \text{ implies that } h=k
        \]
        and we say $f$ is \textbf{monic}
        \item a \textbf{epimorphism} (or just a \textbf{epi}) if for any parallel morphisms $h,k:y \rightrightarrows z$
        \[
            hf=kf \text{ implies that } h=k
        \]
        and we say $f$ is \textbf{epic}
    \end{enumerate}
\end{dfn}

\begin{dfn}[alt] Let \(\mathsf{C}\) be a category with morphism \(f:x \to y\).  Then $f$ is...
    \begin{enumerate}
        \item a \textbf{monomorphism} iff post-composition
        \(
            \left(f_*\right)
        \)
        is an injection on .
        \item a \textbf{epimorphism} iff pre-composition 
        \(\left(f_*\right)\) is injective
    \end{enumerate}
\end{dfn}

\begin{tabular}{|c|c|c|}
    \hline
    &&\\
    mono    & \(\forall h,k:w \rightrightarrows x \quad fh=fk \leadsto h=k \) & \(\rightarrowtail\)\\
    % \(\forall z \quad \forall g:x\to z \quad \exists h:t\to z \st h\circ f=g\) \\
    &&\\
    \hline
    &&\\
    epi     & \(\forall h,k:y \rightrightarrows z \quad hf=kf \leadsto h=k \) & \(\twoheadrightarrow\)\\
    % \(\forall z \quad \forall g:z \to y \quad \exists h:z \to x \st f\circ h =g \) \\
    &&\\
\hline
\end{tabular}

\begin{dfn} suppose \(x \overset{s}{\longrightarrow}y \overset{r}{\longrightarrow} x \st rs=1_x\)
\begin{itemize}
    \item \(s\) is a \textbf{section} or \textbf{right inverse} to \(r\)
    \item \(r\) is a \textbf{retraction} or \textbf{left inverse} to \(s\)
    \item and we say \(s\) and \(r\) express the object x as a \textbf{retract} of the object \(y\)
\end{itemize}
Further, we call \(s\) a \textbf{split monomorphism} and \(r\) a \textbf{split epimorphism}
\end{dfn}

\begin{lem}[1.2.11]\phantom{a}
\begin{itemize}
    \item if \(f,g\) are mono (\(\mono\)) then so is \(g\circ f\)
    \item if \(f,g\) are epi (\(\epi\)) then so is \(g\circ f\)
\end{itemize}
\textit{contrapositively}
\begin{itemize}
    \item if \(g\circ f\) is mono (\(\mono\)) then so is \(f\)
    \item if \(g\circ f\) is epi (\(\epi\)) then so is \(g\)
\end{itemize}
\end{lem}

\begin{prm}[1.2.i]
    Defining \(\mathsf{C}/c\) to be \(\left(c/\left(\mathsf{C}^\op\right)\right)^\op\), deduce Exercise 1.1.iii(ii) from Exercise 1.1.iii(i)
\end{prm}

\begin{prm}[1.2.ii]
    \begin{enumerate}
        \item Show that a morphism \(f:x\to y \) is a split epimorphism in a category \(\mathsf{C}\) if and only if for all \(c\in \mathsf{C}\), the post composition function \(f_*: \mathsf{C}\left(c,x\right)\to \mathsf{C}\left(c,y\right)\) is surjective.
        \item Argue by duality that \(f\) is a split monomorphism if and only if for all \(c\in \mathsf{C}\), the precomposition function \(f^*: \mathsf{C}\left(y,c\right) \to \mathsf{C} \left(x,c\right)\) is surjective
    \end{enumerate}
\end{prm}

\begin{prm}[1.2.iii] Prove Lemma 1.2.11 by proving either either epi or mono pairs of statements then using duality to prove the other pair.  Prove that epi form a subcategory and dually mono
\end{prm}

\begin{prm}
    iv, v, vi, vii
\end{prm}

\section*{1.3 Functoriality}

\begin{dfn} A \textbf{functor} (covariant) \(F: \mathsf{C} \to \mathsf{D}\), between categories $\mathsf{C}$ and \(\mathsf{D}\), consists of the following data:
\begin{enumerate}
    \item[] An object \(Fc\in \mathsf{D}\), for each object \(c\in \mathsf{C}\)
    \item[] A morphism \(Ff:Fc\to Fc' \in \mathsf{D}\), for each morphism \(f:c\to c' \in \mathsf{C}\), so that the domain and codomain of \(Ff\) are, respectively, equal to \(F\) applied to the domain or codomain of \(f\)
\end{enumerate}
The assignments are required to satisfy the following two \textbf{functoriality axioms}
\begin{enumerate}
    \item[] For any composable pair \(f,g \in \mathsf{C}\), \(Fg\cdot Ff=F \left(g\cdot f\right)\)
    \item[] For each object \(c\in \mathsf{C}\), \(F \left(1_c\right)=1_{F_c}\)
\end{enumerate}
\end{dfn}

\begin{dfn} A \textbf{functor} (contravariant) \(F\) from \(\mathsf{C}\) to \(\mathsf{D}\) is a functor \(F: \mathsf{C}^{op} \to \mathsf{D}\). Explicitly, this consists of the following data:
\begin{enumerate}
    \item[] An object \(Fc\in \mathsf{D}\), for each object \(c\in \mathsf{C}\)
    \item[] A morphism \(Ff:Fc'\to Fc \in \mathsf{D}\), for each morphism \(f:c\to c' \in \mathsf{C}\), so that the domain and codomain of \(Ff\) are, respectively, equal to \(F\) applied to the codomain or domain of \(f\)
\end{enumerate}
The assignments are required to satisfy the following two \textbf{functoriality axioms}
\begin{enumerate}
    \item[] For any composable pair \(f,g \in \mathsf{C}\), \(Ff\cdot Fg=F \left(g\cdot f\right)\)
    \item[] For each object \(c\in \mathsf{C}\), \(F \left(1_c\right)=1_{F_c}\)
\end{enumerate}
\end{dfn}

\begin{rmk}
    even in the presence of opposite categories, we always make an effort to draw arrows pointing in the “correct way” and depict composition in the usual order.
\end{rmk}

\begin{lem} Functors preserve isomorphisms.  or
    \[
        \text{if } f:c \to c' \text{ is an isomorphism \underline{then} }  Ff:Fc \to Fc' \text{ is an isomorphism}
    \]
\end{lem}

\begin{pf}
    \item[] Consider a functor \(F:\mathsf{C}\to \mathsf{D}\) and an isomorphism \(f:x \to y\) in \(\mathsf{C}\) with inverse \(g:y\to x\).
    \item[] Apply the two functoriality axioms:
    \[
        F(g)F(f)=F(gf)=F(1_x)=1_{F_x}
    \]
    So, \(fg:Fy \to Fx\) is a left inverse to \(Ff:Fx \to Fy\).  
    \item[] Exchanging the roles of \(f\) and \(g\) shows that \(Fg\) is also a right inverse.
\end{pf}

\begin{dfn} If \(\mathsf{C}\) is locally small, then for any object \(c \in \mathsf{C}\) we define a pair of covariant and contravariant \textbf{functors represented by} \(c\)
\[
\begin{tikzcd}[column sep=large]
\mathcal{C} \arrow[r, "{\mathcal{C}(c,-)}"] & \mathbf{Set}
\end{tikzcd}
\qquad
\begin{tikzcd}[column sep=large]
\mathcal{C}^{\mathrm{op}} \arrow[r, "{\mathcal{C}(-,c)}"] & \mathbf{Set}
\end{tikzcd}
\]
    \[
\begin{tikzcd}[row sep=large, column sep=large]
x & y \\
{\mathcal{C}(c,x)} & {\mathcal{C}(c,y)}
\arrow[mapsto, from=1-1, to=2-1]
\arrow[mapsto, from=1-2, to=2-2]
\arrow["{f}"{name=F}, from=1-1, to=1-2]
\arrow["{f_\ast}"'{name=G}, from=2-1, to=2-2]
\arrow[mapsto, from=F, to=G, shorten <=.6em, shorten >=.6em]
\end{tikzcd}
\qquad
\begin{tikzcd}[row sep=large, column sep=large]
x & y \\
{\mathcal{C}(x,c)} & {\mathcal{C}(y,c)}
\arrow[mapsto, from=1-1, to=2-1]
\arrow[mapsto, from=1-2, to=2-2]
\arrow["{f}"{name=F}, from=1-1, to=1-2]
\arrow["{f^\ast}"{name=H}, from=2-2, to=2-1, swap]
\arrow[mapsto, from=F, to=H, shorten <=.6em, shorten >=.6em]
\end{tikzcd}
\]

\end{dfn}

\begin{prp}  the covariant and contravariant \textbf{functors represented by \(c\)} satisfy the functoriality axioms
    
\end{prp}

\begin{pf}
    (left to me) hint on pg $19$
\end{pf}

\begin{dfn} a functor of two varibles is called a bifunctor
\end{dfn}

\begin{dfn} For any categories \(\mathsf{C}\) and \(\mathsf{D}\), there is a category \(\mathsf{C}\times\mathsf{D}\), their \textbf{product}, whose 
\begin{itemize}
    \item objects are ordered pairs \(\left(c,d\right)\), where \(c\) is an object of \(\mathsf{C}\) and \(d\) is an object of \(\mathsf{D}\)
    \item morphisms are ordered pairs \(\left(f,g\right):\left(c,d\right)\to \left(c',d'\right),\) where \(f:c\to c' \in \mathsf{C}\) and \(g:d\to d' \in \mathsf{D}\)
    \item in which composition and identities are defined componentwise
\end{itemize}
\end{dfn}

\begin{dfn} If \(\mathsf{C}\) is locally small, then there is a \textbf{two-sided represented functor}
    \[
        \mathsf{C} \left(-,-\right): \mathsf{C}^{op}\times \mathsf{C} \to \mathsf{Set}
    \]
    defined in the evident manner.  A pair of objects \(\left(x,y\right)\) is mapped to the hom-set \(\mathsf{C} \left(x,y\right)\).  A pair of morphisms \(f:w\to x\) and \(h:y\to z\) is sent to the function
    \[\begin{tikzcd}
        \mathsf{C}\left(x,y\right) \arrow[r,"{f^*,h_*}"] & \mathsf{C}\left(w,z\right) \\
        g \ar[r,mapsto] & hgf
    \end{tikzcd}\]
that takes an arrow \(g:x\to y\) and then pre-composes with \(f\) and post-composes with \(h\) to obtain \(hgf:w\to z\)
\end{dfn}

\section*{1.4 Naturality}

\subsection*{brief discursion on the linear dual}

\begin{dfn} Let \(V\) be a finite-dimension \(\k\)-vector space.  A linear map \(\psi:V\to \k\) is a \textbf{linear functional} on \(V\)
    \begin{itemize}
        \item the \textbf{linear dual} of \(V\) is given:
        \[\begin{aligned}
            V^*\:&\mathrm{Hom}\left(V,\k\right)\\
            =&\left\{\psi \in \mathcal{F} \left(V\to \k\right)\, : \,\psi \text{ is a linear map}\right\}
        \end{aligned}\]
        \item the \textbf{double dual} of \(V\) is given:
        \[\begin{aligned}
            V^{**} \:& \Hom\left(\Hom\left(V,\k\right),\k\right)\\
            =&\left\{\Psi \in \mathcal{F} \left(V^*\to \k\right)\, : \,\Psi \text{ is a linear map}\right\}
        \end{aligned}\]
    \end{itemize}
\end{dfn}

\begin{dfn} Given categories \(\mathsf{C}\) and \(\mathsf{D}\) and functors \(F,G:\mathsf{C}\rightrightarrows \mathsf{D}\), a \textbf{natural transformation} \(\alpha:F \Rightarrow G\) consists of:
    \begin{itemize}
        \item an arrow \(\alpha_c:Fc\to Gc\) in \(\mathsf{D}\) for each object \(c\in \mathsf{C}\), the collection of which define the \textbf{components} of the natural transformation, so that for any morphism \(f:c\to c'\) in \(\mathsf{C}\), the following square of morphisms in \(\mathsf{D}\):
        \[
            \begin{tikzcd}
                Fc \ar[r,"\alpha_c"] \ar[d,"Ff"'] 
                    & Gc \ar[d,"Gf"] \\
                Fc' \ar[r,"\alpha_{c'}"]
                    & || D
            \end{tikzcd}
        \]
    \end{itemize}
     \textbf{commutes}, i.e., has a common composite \(Fc\to Gc'\) in \(\mathsf{D}\)
\begin{itemize}
    \item More succinctly we write describe just the \textit{boundary data}:
    \item \textcolor{red}{add globular diagram}
    \item A \textbf{natural isomorphism} is a natural transformation \(\alpha:F\implies G\) in which every component \(\alpha_c\) is an isomorphism.  In this case, the natural isomorphism may be depicted as \(\alpha: F \cong G\)
\end{itemize}
    \end{dfn}

\end{document}